% =====================================================================
\chapter{Modulo 11 -- Liste Concatenate}
% =====================================================================

\tikzset{
  nodoL/.style={draw,rectangle split,rectangle split horizontal,rectangle split parts=2,
                fill=heapcol,font=\ttfamily\small,minimum height=0.65cm,
                rectangle split part fill={heapcol,heapcol!60!sfondo}},
  nullo/.style={font=\ttfamily\footnotesize,text=black!60},
}

% ---------------------------------------------------------------------
\section{Teoria \& Analogia Chiave}
% ---------------------------------------------------------------------

Da questo modulo in poi costruiamo le \textbf{strutture dati dinamiche}: liste,
pile, code, alberi. All'esame vanno scritte \textbf{da zero}, con nodi allocati con
\cc|new| e collegati da puntatori, senza la STL. Servono tutto il Modulo 3
(puntatori, \cc|new|/\cc|delete|), il Modulo 4 (distruttore e Regola dei Tre) e il
Modulo 5 (template). Se qualcosa ti sembra oscuro, torna lì.

\subsection*{Perché le liste: il limite degli array}

Un array è una fila di case \textbf{attaccate} l'una all'altra (Modulo 3). È
comodissimo per l'\textbf{accesso diretto}: \cc|v[i]| si raggiunge subito, perché il
suo indirizzo si calcola. Ma ha due difetti:
\begin{itemize}
  \item la dimensione è fissa: per farlo crescere bisogna allocare un nuovo blocco e
        ricopiare tutto (lo abbiamo fatto con il raddoppio);
  \item inserire o togliere un elemento \emph{in mezzo} obbliga a spostare tutti quelli
        dopo (Esercizio~3.10): $O(n)$.
\end{itemize}

Una \textbf{lista concatenata} risolve il problema rinunciando alla contiguità. Ogni
elemento, il \textbf{nodo}, vive dove capita nello heap e contiene due cose: il
\textbf{valore} e l'\textbf{indirizzo del nodo successivo}. Basta conoscere
l'indirizzo del primo nodo, la \textbf{testa}, per raggiungere tutti gli altri
seguendo le frecce. L'ultimo nodo ha come successivo \cc|nullptr|: ``qui finisce''.

\begin{analogia}[: la caccia al tesoro]
Un array è una fila di armadietti numerati: vai direttamente al numero 7. Una lista è
una caccia al tesoro: hai solo il primo biglietto (la testa), che contiene un premio
(il valore) e l'indicazione per trovare il biglietto successivo (il puntatore). Per
arrivare al settimo devi seguire i primi sei (accesso \textbf{sequenziale}, $O(i)$).
In compenso, aggiungere un biglietto in mezzo è facilissimo: basta cambiare
l'indicazione scritta su quello prima.
\end{analogia}

\begin{center}\small
\begin{tabular}{p{5.4cm}p{3.6cm}p{5.4cm}}
\toprule
\textbf{Operazione} & \textbf{Array} & \textbf{Lista concatenata} \\
\midrule
accesso all'elemento $i$-esimo & $O(1)$ & $O(i)$ \\
inserimento/rimozione in testa & $O(n)$ (sposto tutti) & $O(1)$ \\
inserimento/rimozione dopo un nodo noto & $O(n)$ & $O(1)$ \\
inserimento in coda & $O(1)$ se c'è posto & $O(n)$, oppure $O(1)$ con puntatore alla coda \\
ricerca di un valore & $O(n)$ ($O(\log n)$ se ordinato) & $O(n)$ (anche se ordinata) \\
memoria & solo i valori & valore + un puntatore per nodo \\
\bottomrule
\end{tabular}
\end{center}
Nota l'ultima riga della ricerca: su una lista ordinata \textbf{non} si può fare la
ricerca binaria, perché non si può saltare a metà.

\subsection*{Il nodo e la lista}

Le slide definiscono il nodo con una \cc|struct| o una \cc|class| con attributi
pubblici: è un semplice ``pacchetto'' di dati, senza regole da proteggere.
\begin{lstlisting}
template <typename T>
struct Nodo {
    T valore;
    Nodo<T>* succ;                                  // indirizzo del successivo
    Nodo(T v, Nodo<T>* s = nullptr) : valore(v), succ(s) {}
};
\end{lstlisting}
Attenzione: \cc|Nodo<T>* succ| è un \emph{puntatore} a un nodo, non un nodo. Un nodo
che contenesse un altro nodo intero sarebbe infinito (e il compilatore lo vieta).

La \textbf{lista} vera e propria è una classe che custodisce il puntatore alla testa
(ed eventualmente alla coda e il numero di nodi) e offre i metodi. Gli utenti della
lista non toccano mai i nodi direttamente: information hiding.
\begin{lstlisting}
template <typename T>
class Lista {
    Nodo<T>* testa;
public:
    Lista() : testa(nullptr) {}     // lista vuota: testa nulla
    ...
};
\end{lstlisting}

\subsection*{Scorrere una lista: il ``for'' delle liste}

Questo è lo schema che userai \textbf{ovunque}. Memorizzalo:
\begin{lstlisting}
for (Nodo<T>* p = testa; p != nullptr; p = p->succ) {
    // qui p punta a un nodo: p->valore è il suo valore
}
\end{lstlisting}
\cc|p| è un ``dito'' che parte dalla testa e a ogni giro salta al nodo successivo, finché
non cade fuori dalla lista (\cc|nullptr|). Non si muove mai \cc|testa|: se lo facessi,
perderesti l'inizio della lista per sempre.

\subsection*{Le operazioni fondamentali, passo per passo}

Per ognuna, nella prossima sezione c'è il disegno. Il segreto delle liste è
\textbf{disegnare prima di scrivere codice}: scatolette con frecce, e quali frecce
cambiano in quale ordine.

\textbf{Inserimento in testa} --- $O(1)$, il più semplice:
\begin{lstlisting}
testa = new Nodo<T>(v, testa);   // il nuovo nodo punta alla vecchia testa e diventa la testa
\end{lstlisting}
Funziona anche con la lista vuota (la vecchia testa è \cc|nullptr|, quindi il nuovo
nodo diventa l'unico).

\textbf{Inserimento dopo un nodo \texttt{p}} --- $O(1)$, l'ordine conta:
\begin{lstlisting}
Nodo<T>* nuovo = new Nodo<T>(v);
nuovo->succ = p->succ;   // 1. il nuovo punta a quello che seguiva p
p->succ = nuovo;         // 2. p punta al nuovo
\end{lstlisting}
Se inverti le due righe, \cc|p->succ| viene sovrascritto prima di essere copiato: perdi
tutto il resto della lista.

\textbf{Inserimento in coda} --- si cammina fino all'\emph{ultimo} nodo (quello con
\cc|succ == nullptr|) e lo si aggancia lì. Attenzione al caso lista vuota. Con un
attributo \cc|coda| che punta sempre all'ultimo nodo diventa $O(1)$, ma va aggiornato
in \emph{tutte} le operazioni che possono cambiare l'ultimo nodo.

\textbf{Inserimento ordinato} --- si cerca il nodo \cc|p| dopo il quale inserire
(l'ultimo con valore minore) e si inserisce dopo \cc|p|. Se il nuovo valore va prima
della testa, è un inserimento in testa.

\textbf{Rimozione di un valore} --- la parte più delicata. Per togliere un nodo devo
far ``scavalcare'' quel nodo dal suo \textbf{precedente}: quindi mentre cerco tengo due
puntatori, \cc|ant| (il precedente) e \cc|pos| (il corrente), come nelle slide:
\begin{passopasso}[: rimuovere il nodo con valore v]
\begin{enumerate}
  \item scorri con \cc|ant| e \cc|pos| finché \cc|pos| non punta al nodo con \cc|v|
        (o diventa \cc|nullptr|: non trovato);
  \item se \cc|ant == nullptr| il nodo da togliere è la testa: \cc|testa = pos->succ;|
  \item altrimenti scavalca: \cc|ant->succ = pos->succ;|
  \item \textbf{solo ora} \cc|delete pos;|
\end{enumerate}
\end{passopasso}

\textbf{Distruzione della lista} --- il distruttore deve liberare \emph{ogni} nodo.
Salva il successivo \textbf{prima} di cancellare il corrente, altrimenti leggeresti
\cc|p->succ| da un nodo già liberato:
\begin{lstlisting}
while (testa != nullptr) {
    Nodo<T>* tmp = testa;
    testa = testa->succ;   // prima avanzo...
    delete tmp;            // ...poi cancello quello che ho lasciato
}
\end{lstlisting}

\textbf{Inversione in loco} --- giro tutte le frecce senza creare nodi, con tre
puntatori \cc|prec|, \cc|corr|, \cc|succ| che avanzano insieme. È una domanda d'esame
classica:
\begin{lstlisting}
Nodo<T>* prec = nullptr;
Nodo<T>* corr = testa;
while (corr != nullptr) {
    Nodo<T>* succ = corr->succ;   // 1. ricordo dove andare
    corr->succ = prec;            // 2. giro la freccia
    prec = corr;                  // 3. avanzo i due puntatori
    corr = succ;
}
testa = prec;
\end{lstlisting}

\begin{consiglio}[: i tre casi da controllare SEMPRE]
Ogni operazione su una lista va provata (a mente o con un \cc|main|) su: \textbf{lista
vuota}, \textbf{un solo nodo}, \textbf{operazione sulla testa} (e sulla coda). Il 90\%
dei bug delle liste sta in uno di questi tre casi.
\end{consiglio}

\subsection*{La lista possiede i suoi nodi: Regola dei Tre}

I nodi sono nello heap e li ha creati la lista: è la lista a doverli distruggere. Quindi
servono distruttore, costruttore di copia (che crea nodi \emph{nuovi} con gli stessi
valori: copia profonda) e \cc|operator=| (Modulo 4). Se non vuoi permettere la copia,
puoi vietarla esplicitamente:
\cc|Lista(const Lista&) = delete;| (lo fa il laboratorio): così un tentativo di copia
è un errore di compilazione invece di un crash.

\subsection*{Varianti}

\begin{itemize}
  \item \textbf{Lista doppiamente concatenata}: ogni nodo ha anche il puntatore al
        \textbf{precedente} (\cc|prec|). Costa un puntatore in più per nodo, ma si può
        scorrere all'indietro e rimuovere un nodo conoscendo solo il suo indirizzo
        (il precedente è \cc|pos->prec|). In ogni inserimento o rimozione vanno
        aggiornate \textbf{quattro} frecce invece di due. È la lista del codice del prof
        (Esercizio~11.7).
  \item \textbf{Lista circolare}: l'ultimo nodo punta di nuovo al primo. Non c'è
        \cc|nullptr|: lo scorrimento si ferma quando si torna al punto di partenza.
        Utile per i problemi ``a giro'' (Esercizio~11.9).
  \item \textbf{Lista con sentinella}: un nodo speciale, sempre presente, che non
        contiene dati e marca inizio e fine. In una lista doppia circolare con
        sentinella il primo elemento è \cc|s->succ| e l'ultimo è \cc|s->prec|, entrambi
        raggiungibili in $O(1)$, e \textbf{scompaiono tutti i casi speciali}
        (lista vuota, testa, coda): ogni nodo ha sempre un precedente e un successivo
        (Esercizio~11.10).
\end{itemize}

\begin{inparole}
\begin{itemize}
  \item Nodo = valore + puntatore al successivo. Lista = puntatore alla testa.
  \item Scorrere: \cc|for (p = testa; p != nullptr; p = p->succ)|. Non muovere mai \cc|testa|.
  \item In testa: \cc|testa = new Nodo(v, testa);|. Dopo \cc|p|: prima
        \cc|nuovo->succ = p->succ|, poi \cc|p->succ = nuovo|.
  \item Rimuovere: serve il precedente (\cc|ant|); caso speciale se è la testa;
        \cc|delete| per ultimo.
  \item Distruggere: salva il successivo, poi \cc|delete|.
  \item Prova sempre: lista vuota, un nodo, testa/coda.
\end{itemize}
\end{inparole}

% ---------------------------------------------------------------------
\section{Rappresentazione Visiva della Memoria (TikZ)}
% ---------------------------------------------------------------------

\textbf{Una lista di tre nodi.} Ogni nodo ha una parte \emph{valore} e una parte
\emph{succ}; la freccia parte dalla parte succ. I nodi stanno nello heap in posizioni
qualunque, non contigue.
\begin{center}
\begin{tikzpicture}
  \node[cella,fill=stackcol] (t) at (0,0) {testa};
  \node[nodoL] (a) at (2.6,0) {1 \nodepart{two} \phantom{x}};
  \node[nodoL] (b) at (5.0,0) {3 \nodepart{two} \phantom{x}};
  \node[nodoL] (c) at (7.4,0) {7 \nodepart{two} \phantom{x}};
  \node[nullo] (z) at (9.4,0) {nullptr};
  \draw[freccia] (t) -- (a);
  \draw[freccia] (a.two) -- (b.west);
  \draw[freccia] (b.two) -- (c.west);
  \draw[freccia] (c.two) -- (z.west);
  \node[etichetta,below=0.15cm of a] {valore | succ};
\end{tikzpicture}
\end{center}

\needspace{6cm}
\textbf{Inserimento dopo un nodo \texttt{p}} (inserisco 5 dopo il 3). Le frecce
numerate vanno impostate \textbf{in quell'ordine}.
\begin{center}
\begin{tikzpicture}
  \node[nodoL] (a) at (0,0) {1 \nodepart{two} \phantom{x}};
  \node[nodoL] (b) at (2.6,0) {3 \nodepart{two} \phantom{x}};
  \node[nodoL] (c) at (7.8,0) {7 \nodepart{two} \phantom{x}};
  \node[nodoL,fill=green!25!sfondo,rectangle split part fill={green!25!sfondo,green!15!sfondo}] (n) at (5.2,-1.4) {5 \nodepart{two} \phantom{x}};
  \node[etichetta,above=0.1cm of b] {p};
  \node[etichetta,below=0.1cm of n] {nuovo};
  \draw[freccia] (a.two) -- (b.west);
  \draw[freccia,dashed,black!40] (b.two) -- node[above,etichetta]{vecchia freccia (sostituita)} (c.west);
  \draw[freccia,green!50!black] (n.two) to[out=0,in=-110] node[right,etichetta,text=green!45!black,pos=0.6]{1. nuovo->succ = p->succ} (c.south);
  \draw[freccia,green!50!black] (b.two) to[out=-60,in=180] node[left,etichetta,text=green!45!black,pos=0.4]{2. p->succ = nuovo} (n.west);
\end{tikzpicture}
\end{center}

\needspace{6cm}
\textbf{Rimozione del nodo \texttt{pos}} (tolgo il 3): \cc|ant| lo scavalca, poi lo si
cancella.
\begin{center}
\begin{tikzpicture}
  \node[nodoL] (a) at (0,0) {1 \nodepart{two} \phantom{x}};
  \node[nodoL,fill=red!20!sfondo,rectangle split part fill={red!20!sfondo,red!12!sfondo}] (b) at (2.8,0) {3 \nodepart{two} \phantom{x}};
  \node[nodoL] (c) at (5.6,0) {7 \nodepart{two} \phantom{x}};
  \node[etichetta,above=0.1cm of a] {ant};
  \node[etichetta,above=0.1cm of b] {pos};
  \draw[freccia,dashed,black!40] (a.two) -- (b.west);
  \draw[freccia] (b.two) -- (c.west);
  \draw[freccia,green!50!black] (a.two) to[out=-50,in=-130] node[below,etichetta,text=green!45!black]{1. ant->succ = pos->succ} (c.south);
  \node[etichetta,text=red!75!black,anchor=west] at (7,0) {2. delete pos};
\end{tikzpicture}
\end{center}

\needspace{7cm}
\textbf{Inversione in loco}: tre momenti del ciclo. Le frecce già girate puntano a
sinistra.
\begin{center}
\begin{tikzpicture}
  \foreach \r/\lab in {0/{inizio},1/{dopo 1 giro},2/{fine}} \node[etichetta,anchor=east] at (-0.8,-\r*1.3) {\lab};
  % riga 0
  \node[nodoL] (a0) at (0.6,0) {1 \nodepart{two} \phantom{x}};
  \node[nodoL] (b0) at (3,0) {2 \nodepart{two} \phantom{x}};
  \node[nodoL] (c0) at (5.4,0) {3 \nodepart{two} \phantom{x}};
  \draw[freccia] (a0.two) -- (b0.west); \draw[freccia] (b0.two) -- (c0.west);
  \node[nullo,right=0.4cm of c0] (z0) {nullptr}; \draw[freccia] (c0.two) -- (z0);
  \node[etichetta,above=0.05cm of a0] {corr};
  % riga 1
  \node[nodoL] (a1) at (0.6,-1.3) {1 \nodepart{two} \phantom{x}};
  \node[nodoL] (b1) at (3,-1.3) {2 \nodepart{two} \phantom{x}};
  \node[nodoL] (c1) at (5.4,-1.3) {3 \nodepart{two} \phantom{x}};
  \node[nullo] (z1) at (-0.6,-1.95) {nullptr};
  \draw[freccia,green!50!black] (a1.two) to[out=-90,in=0] (z1);
  \draw[freccia] (b1.two) -- (c1.west);
  \node[etichetta,above=0.05cm of a1] {prec};
  \node[etichetta,above=0.05cm of b1] {corr};
  % riga 2
  \node[nodoL] (a2) at (0.6,-2.6) {1 \nodepart{two} \phantom{x}};
  \node[nodoL] (b2) at (3,-2.6) {2 \nodepart{two} \phantom{x}};
  \node[nodoL] (c2) at (5.4,-2.6) {3 \nodepart{two} \phantom{x}};
  \node[nullo] (z2) at (-0.6,-3.25) {nullptr};
  \draw[freccia,green!50!black] (a2.two) to[out=-90,in=0] (z2);
  \draw[freccia,green!50!black] (b2.two) to[out=-110,in=-60] (a2.south);
  \draw[freccia,green!50!black] (c2.two) to[out=-110,in=-60] (b2.south);
  \node[etichetta,above=0.05cm of c2] {testa = prec};
\end{tikzpicture}
\end{center}

\needspace{5cm}
\textbf{Lista doppia e lista circolare con sentinella.}
\begin{center}
\begin{tikzpicture}
  % doppia
  \node[nullo] (zn) at (-1.2,0) {nullptr};
  \foreach \v/\x [count=\i] in {a/0.6,b/2.6,c/4.6}
    \node[cella,fill=heapcol,minimum width=1.2cm] (d\i) at (\x,0) {\v};
  \node[nullo] (zf) at (6.4,0) {nullptr};
  \draw[freccia] ($(d1.east)+(0,0.12)$) -- ($(d2.west)+(0,0.12)$);
  \draw[freccia] ($(d2.west)+(0,-0.12)$) -- ($(d1.east)+(0,-0.12)$);
  \draw[freccia] ($(d2.east)+(0,0.12)$) -- ($(d3.west)+(0,0.12)$);
  \draw[freccia] ($(d3.west)+(0,-0.12)$) -- ($(d2.east)+(0,-0.12)$);
  \draw[freccia] (d1.west) -- (zn); \draw[freccia] (d3.east) -- (zf);
  \node[etichetta,above=0.1cm of d2] {doppia: succ e prec};
  % circolare con sentinella
  \node[cella,fill=black!15!sfondo,minimum width=1.2cm] (s) at (9.6,0.6) {S};
  \node[cella,fill=heapcol,minimum width=1.2cm] (e1) at (11.4,-0.6) {a};
  \node[cella,fill=heapcol,minimum width=1.2cm] (e2) at (7.8,-0.6) {b};
  \draw[freccia] (s) to[bend left=20] (e1);
  \draw[freccia] (e1) to[bend left=20] (e2);
  \draw[freccia] (e2) to[bend left=20] (s);
  \node[etichetta,above=0.1cm of s] {sentinella: primo = S->succ, ultimo = S->prec};
\end{tikzpicture}
\end{center}

% ---------------------------------------------------------------------
\section{Pattern di Codice di Laboratorio}
% ---------------------------------------------------------------------

La lista semplice generica completa: \cc|Nodo<T>| con costruttore, classe \cc|Lista<T>|
con Regola dei Tre, inserimento in testa, in coda e ordinato, ricerca, rimozione con
\cc|ant|/\cc|pos|, inversione in loco e \cc|operator<<| \cc|friend| scritto dentro la
classe (la forma più semplice con i template).

\codice{codice/m11/pattern.cpp}
\uscita{codice/m11/output/pattern.txt}

\textbf{Cosa osservare:}
\begin{itemize}
  \item \cc|rimuovi(42)| restituisce \cc|false| e non fa danni: il valore non c'è.
  \item Dopo \cc|l.inverti()| la copia è rimasta \texttt{[3 -> 5 -> 7]}: il costruttore
        di copia ha creato nodi nuovi.
  \item \cc|operator=| usa un trucco elegante: crea una copia, scambia le teste e lascia
        che il distruttore della copia liberi i vecchi nodi. Anche la classica versione a
        sei passi del Modulo 4 va benissimo.
  \item Lo stesso codice funziona con \cc|string|: \cc|inserisciOrdinato| usa solo
        \cc|<|.
\end{itemize}

\begin{consiglio}[: operator<< con i template, due modi]
Il modo più semplice (usato qui) è scrivere la funzione \cc|friend| \emph{dentro} la
classe. Il codice del prof la dichiara invece come template separato:
\cc|template <typename U> friend ostream& operator<<(ostream&, const ListaDoppiaT<U>&);|
e la definisce fuori con \cc|template <typename U>|. Funzionano entrambi.
\end{consiglio}

% ---------------------------------------------------------------------
\section{Tips \& Trappole da Esame}
% ---------------------------------------------------------------------

\begin{trappola}[: usare un nodo dopo il delete]
\cc|delete p; p = p->succ;| legge da memoria già liberata: comportamento indefinito,
spesso un crash solo ``ogni tanto''. Salva sempre \cc|p->succ| \emph{prima} del
\cc|delete|.
\end{trappola}

\begin{trappola}[: perdere la lista]
\cc|p->succ = nuovo;| prima di \cc|nuovo->succ = p->succ;| fa perdere tutti i nodi dopo
\cc|p| (memory leak, e la lista è troncata). L'ordine delle assegnazioni conta:
disegna.
\end{trappola}

\begin{trappola}[: muovere la testa per scorrere]
\cc|while (testa) { cout << testa->valore; testa = testa->succ; }| stampa la lista e
poi... la lista è persa: \cc|testa| è \cc|nullptr|. Per scorrere usa un puntatore
ausiliario \cc|p|.
\end{trappola}

\begin{trappola}[: lista vuota e primo nodo]
\cc|testa->succ| con \cc|testa == nullptr| è un crash. Le operazioni in coda, la
rimozione e l'inserimento ordinato hanno quasi sempre un caso speciale per la lista
vuota e per la testa.
\end{trappola}

\begin{trappola}[: la funzione che cambia la testa]
Una funzione esterna \cc|void inserisciTesta(Nodo* testa, int v)| modifica solo la sua
copia del puntatore: il chiamante non vede nulla. Serve \cc|Nodo*& testa| (Modulo 3) o
restituire la nuova testa. Dentro una classe il problema non c'è, perché \cc|testa| è
un attributo.
\end{trappola}

\begin{trappola}[: distruttore mancante]
Senza distruttore, alla fine dello scope tutti i nodi restano nello heap. E con il
distruttore ma senza costruttore di copia, copiare una lista porta a un doppio
\cc|delete|. Regola dei Tre, oppure \cc|= delete| sulla copia.
\end{trappola}

\begin{consiglio}[: stampa la lista dopo ogni operazione]
Mentre sviluppi, stampa la lista dopo \emph{ogni} inserimento e rimozione: un errore
nei puntatori si vede subito (nodi spariti, ripetuti, ciclo infinito nella stampa).
\end{consiglio}

% ---------------------------------------------------------------------
\section{Quiz a Crocette (Prova Scritta)}
% ---------------------------------------------------------------------

\begin{quiz}{Qual è il costo dell'accesso all'$i$-esimo elemento in una lista concatenata semplice?}
  \item $O(1)$
  \item $O(\log i)$
  \item $O(i)$
  \item $O(i^2)$
\end{quiz}

\begin{quiz}{Qual è il costo dell'inserimento in testa in una lista concatenata?}
  \item $O(1)$
  \item $O(\log n)$
  \item $O(n)$
  \item dipende dal valore inserito
\end{quiz}

\begin{quiz}{Si può applicare la ricerca binaria a una lista concatenata ordinata ottenendo $O(\log n)$?}
  \item sì, come con gli array
  \item no, perché non si può accedere in tempo costante all'elemento centrale
  \item sì, ma solo con le liste doppie
  \item no, perché le liste non possono essere ordinate
\end{quiz}

\begin{quiz}{Per inserire un nodo \texttt{nuovo} dopo il nodo \texttt{p}, quale sequenza è corretta?}
  \item \texttt{p->succ = nuovo; nuovo->succ = p->succ;}
  \item \texttt{nuovo->succ = p->succ; p->succ = nuovo;}
  \item \texttt{nuovo->succ = p; p->succ = nuovo;}
  \item \texttt{p = nuovo; nuovo->succ = p;}
\end{quiz}

\begin{quiz}{Per rimuovere un nodo da una lista concatenata semplice, di quale nodo serve conoscere l'indirizzo, oltre a quello da rimuovere?}
  \item del successivo
  \item del precedente
  \item dell'ultimo
  \item di nessuno
\end{quiz}

\begin{quiz}{Cosa c'è di sbagliato in \texttt{while (p) \{ delete p; p = p->succ; \}}?}
  \item niente
  \item si legge \texttt{p->succ} da un nodo già liberato
  \item non libera l'ultimo nodo
  \item \texttt{delete} non si può usare in un ciclo
\end{quiz}

\begin{quiz}{Quale vantaggio ha una lista doppiamente concatenata rispetto a una semplice?}
  \item occupa meno memoria
  \item permette l'accesso diretto all'$i$-esimo elemento
  \item permette di scorrere la lista in entrambe le direzioni e di rimuovere un nodo noto senza cercare il precedente
  \item non richiede il puntatore alla testa
\end{quiz}

\begin{quiz}{In una lista circolare semplice, come si capisce di aver finito di scorrerla?}
  \item quando si trova \texttt{nullptr}
  \item quando si torna al nodo di partenza
  \item quando il valore è zero
  \item non si può scorrere
\end{quiz}

\begin{quiz}{A cosa serve il nodo sentinella in una lista doppia circolare?}
  \item a contenere il numero di elementi
  \item a eliminare i casi speciali (lista vuota, testa, coda) e ad accedere a primo e ultimo in $O(1)$
  \item a rendere la lista ordinata
  \item a impedire le copie della lista
\end{quiz}

\begin{quiz}{Qual è la complessità dell'inversione in loco di una lista di $n$ nodi?}
  \item $O(1)$
  \item $O(\log n)$
  \item $O(n)$ in tempo e $O(1)$ di memoria aggiuntiva
  \item $O(n)$ in tempo e $O(n)$ di memoria aggiuntiva
\end{quiz}

\begin{soluzione}[: risposte commentate]
\begin{enumerate}[label=\textbf{\arabic*.}]
  \item \textbf{C}. Accesso sequenziale: bisogna seguire $i$ frecce.
  \item \textbf{A}. Un nodo nuovo e due assegnazioni.
  \item \textbf{B}. Raggiungere il centro costa $O(n)$.
  \item \textbf{B}. Prima si copia \cc|p->succ|, poi lo si sovrascrive.
  \item \textbf{B}. Il precedente deve scavalcare il nodo rimosso.
  \item \textbf{B}. Va salvato il successivo prima del \cc|delete|.
  \item \textbf{C}. Costa invece \emph{più} memoria (A falsa).
  \item \textbf{B}. Non c'è nessun \cc|nullptr|.
  \item \textbf{B}.
  \item \textbf{C}. Tre puntatori ausiliari, indipendenti da $n$.
\end{enumerate}
\end{soluzione}

% ---------------------------------------------------------------------
\section{Laboratorio a Step (Simulazione d'Esame)}
% ---------------------------------------------------------------------

\begin{traccia}[: rettangoli in una lista (dalle slide)]
Definire una semplice classe \cc|Rettangolo| e una lista concatenata generica scritta da
zero. Definire un metodo statico che legge da input un numero arbitrario di rettangoli e
li inserisce nella lista, e un metodo statico che elimina dalla lista tutti i rettangoli
con area compresa tra due valori dati.
\end{traccia}

\begin{step}{la lista minima}
\cc|template <typename T> struct Nodo| e \cc|class Lista| con \cc|inserisciTesta|,
\cc|operator<<| e distruttore. Vieta la copia con \cc|= delete|. Provala con degli
interi.\\[0.3em]
\textbf{Concetto nuovo:} nodo, testa, scorrimento, distruzione.
\end{step}

\begin{step}{la coda}
Aggiungi un attributo \cc|coda| (puntatore all'ultimo nodo), \cc|inserisciCoda| in
$O(1)$ e \cc|size()|. Aggiorna \cc|coda| anche in \cc|inserisciTesta| quando la lista era
vuota.\\[0.3em]
\textbf{Concetto nuovo:} mantenere più puntatori coerenti.
\end{step}

\begin{step}{rimozione multipla}
Scrivi \cc|int rimuoviSe(bool (*condizione)(const T&, double, double), double a, double b)|,
che toglie \textbf{tutti} i nodi per cui la condizione è vera e restituisce quanti ne ha
tolti. Attenzione: dopo una rimozione \cc|ant| non avanza; e se togli l'ultimo nodo,
aggiorna \cc|coda|.\\[0.3em]
\textbf{Concetto nuovo:} rimozione durante lo scorrimento, puntatore a funzione.
\end{step}

\begin{step}{i rettangoli}
Scrivi \cc|Rettangolo| (con \cc|area()| e \cc|operator<<|) e una classe \cc|Gestore| con i
metodi statici \cc|leggi(Lista<Rettangolo>&)| (legge il numero di rettangoli e poi le
coppie base altezza) ed \cc|eliminaTra(Lista<Rettangolo>&, double a, double b)|.\\[0.3em]
\textbf{Concetto nuovo:} lista di oggetti, metodi statici, la traccia delle slide.
\end{step}

\begin{step}{lista ordinata}
Aggiungi \cc|operator<| a \cc|Rettangolo| (per area) e \cc|inserisciOrdinato| alla lista,
e costruisci una lista di rettangoli ordinata per area.\\[0.3em]
\textbf{Concetto nuovo:} inserimento ordinato.
\end{step}

\begin{step}{inversione}
Scrivi \cc|inverti()| in loco, aggiornando correttamente \cc|coda| (la vecchia testa).
Verifica con un \cc|inserisciCoda| dopo l'inversione.\\[0.3em]
\textbf{Concetto nuovo:} manipolazione di tutti i puntatori.
\end{step}

\subsection*{Soluzione completa}
\codice{codice/m11/lab.cpp}
\ingresso{codice/m11/input/lab.txt}
\uscita{codice/m11/output/lab.txt}

\begin{consiglio}[: cosa notare nella soluzione]
\begin{itemize}
  \item In \cc|rimuoviSe| il puntatore \cc|pos| avanza \emph{prima} del \cc|delete|,
        e \cc|ant| avanza solo quando il nodo resta.
  \item Il test finale dello step 6 (un inserimento in coda dopo l'inversione) serve
        proprio a scoprire se \cc|coda| è stata aggiornata: se puntasse ancora al vecchio
        ultimo nodo, il nuovo rettangolo finirebbe nel posto sbagliato.
  \item Il \cc|Gestore| non ha attributi: è solo un contenitore di metodi statici, come
        chiede la traccia.
\end{itemize}
\end{consiglio}

% ---------------------------------------------------------------------
\section{Domande e Risposte Modello per l'Orale}
% ---------------------------------------------------------------------

\begin{orale}[D1. Differenze tra array e lista concatenata?]
Nell'array gli elementi occupano locazioni contigue, quindi l'accesso all'$i$-esimo è
diretto, in $O(1)$; la dimensione però è fissa, e inserimenti e rimozioni in mezzo
richiedono di spostare elementi, in $O(n)$. Nella lista ogni nodo contiene il valore e il
puntatore al successivo, e i nodi possono stare ovunque in memoria: l'accesso è
sequenziale, in $O(i)$, ma la dimensione cresce e cala dinamicamente e
inserimenti e rimozioni, noto il punto in cui farli, costano $O(1)$. La lista occupa più
memoria per via dei puntatori e non permette la ricerca binaria.
\end{orale}

\begin{orale}[D2. Come si rimuove un elemento da una lista semplice?]
Si scorre la lista con due puntatori: \cc|pos| sul nodo corrente e \cc|ant| sul
precedente, finché \cc|pos| non punta al nodo da rimuovere. Se \cc|ant| è nullo il nodo
è la testa, e la testa diventa \cc|pos->succ|; altrimenti si fa puntare \cc|ant->succ| a
\cc|pos->succ|, scavalcando il nodo. Infine si libera la memoria del nodo con
\cc|delete|. La ricerca costa $O(n)$; lo scollegamento $O(1)$.
\end{orale}

\begin{orale}[D3. Cos'è una lista doppiamente concatenata? Vantaggi e svantaggi?]
È una lista in cui ogni nodo contiene, oltre al valore, due puntatori: al successivo e
al precedente. Permette di scorrere la lista in entrambe le direzioni e di rimuovere un
nodo di cui si conosce l'indirizzo in $O(1)$, senza cercarne il predecessore. Gli
svantaggi sono la memoria per il puntatore aggiuntivo e il fatto che ogni inserimento o
rimozione richiede di aggiornare più collegamenti.
\end{orale}

\begin{orale}[D4. Cos'è una lista circolare con sentinella?]
In una lista circolare l'ultimo nodo è collegato al primo. La sentinella è un nodo
speciale, sempre presente, che non contiene informazione utile e marca l'inizio e la fine
della lista. In una lista doppia circolare con sentinella il primo elemento è il
successore della sentinella e l'ultimo il suo predecessore, entrambi accessibili in
$O(1)$; inoltre ogni nodo ha sempre un precedente e un successivo, per cui spariscono i
casi particolari di lista vuota, inserimento in testa e in coda.
\end{orale}

\begin{orale}[D5. Perché una classe lista ha bisogno della Regola dei Tre?]
Perché possiede risorse dinamiche: i nodi allocati con \cc|new|. Il distruttore deve
scorrere la lista e liberare ogni nodo. Senza un costruttore di copia e un operatore di
assegnamento scritti a mano, la copia generata dal compilatore copierebbe solo il
puntatore alla testa: le due liste condividerebbero i nodi e alla distruzione si avrebbe
una doppia deallocazione. La copia corretta crea nuovi nodi con gli stessi valori; in
alternativa si può vietare la copia dichiarando quei metodi \cc|= delete|.
\end{orale}

% ---------------------------------------------------------------------
\section{Esercizi al PC}
% ---------------------------------------------------------------------

Per ogni esercizio, \textbf{prima} disegna su carta i nodi e le frecce che cambiano.

\begin{esercizio}{La prima lista, a mano}{facile}
Senza classi: dichiara \cc|struct Nodo { int valore; Nodo* succ; };|, costruisci la lista
$4 \to 3 \to 5$ inserendo in testa (come nelle slide), aggiungi 8 in coda, scorrila
contando i nodi e sommando i valori, poi libera tutti i nodi.
\end{esercizio}

\begin{esercizio}{Ricerca e modifica}{facile}
Scrivi \cc|Nodo* ricerca(Nodo* testa, int x)| e \cc|int posizione(Nodo* testa, int x)|.
Usa il puntatore restituito da \cc|ricerca| per modificare il valore del nodo
direttamente nella lista.
\end{esercizio}

\begin{esercizio}{Massimo}{facile}
Scrivi il massimo di una lista non vuota in versione iterativa e in versione ricorsiva
(caso base: un solo nodo).
\end{esercizio}

\begin{esercizio}{Togliere in testa e in coda}{facile}
Scrivi \cc|bool rimuoviTesta(Nodo*& testa)| e \cc|bool rimuoviCoda(Nodo*& testa)| (che
si ferma sul \emph{penultimo} nodo). Perché la testa va passata per reference? Svuota una
lista a forza di \cc|rimuoviCoda|.
\end{esercizio}

\begin{esercizio}{Lista ordinata senza duplicati}{medio}
Classe template \cc|ListaOrdinata<T>| (con il nodo come \cc|struct| interna) con
\cc|bool inserisci(const T&)| che mantiene l'ordine e rifiuta i duplicati, e
\cc|bool rimuovi(const T&)| che sfrutta l'ordinamento per fermarsi prima.
\end{esercizio}

\begin{esercizio}{Invertire}{medio}
Inverti una lista in loco in due modi: iterativo con tre puntatori e ricorsivo
(inverti il resto, poi aggancia il nodo corrente in fondo). Le due funzioni
restituiscono la nuova testa.
\end{esercizio}

\begin{esercizio}{Lista doppia del prof}{medio}
Completa la \cc|ListaDoppiaT| del codice del prof: aggiungi \cc|inserisci(T valPrec, T val)|,
che inserisce \cc|val| come \textbf{precedente} del nodo che contiene \cc|valPrec| (se
esiste), una rimozione che sfrutta \cc|prec|, e la stampa in avanti e all'indietro
(tieni un puntatore alla coda).
\end{esercizio}

\begin{esercizio}{Fusione di liste ordinate}{medio}
Fondi due liste ordinate in una lista ordinata \textbf{senza creare nodi}, ricollegando
quelli esistenti. Usa un nodo ``finto'' locale come appoggio per evitare il caso speciale
della testa.
\end{esercizio}

\begin{esercizio}{Il problema di Giuseppe}{difficile}
$n$ persone sono in cerchio, numerate da 1 a $n$. Partendo dalla prima, si conta fino a
$k$ e la persona contata esce; si riprende a contare dalla successiva. Con una lista
\textbf{circolare} simula il gioco, stampa l'ordine di uscita e trova il sopravvissuto per
$n = 7, k = 3$ e $n = 10, k = 2$.
\end{esercizio}

\begin{esercizio}{Lista con sentinella}{difficile}
Classe template di lista doppia circolare con sentinella: costruttore (la sentinella
punta a se stessa), \cc|inserisciTesta| e \cc|inserisciCoda| in $O(1)$ scritti usando
un'unica funzione privata \cc|inserisciDopo|, \cc|rimuovi|, distruttore e \cc|<<| che
stampa in avanti e all'indietro.
\end{esercizio}

\subsection*{Soluzioni}
Le soluzioni commentate, con il codice e l'output reale, sono nel PDF separato
\textbf{Soluzioni} (\texttt{soluzioni-light.pdf} / \texttt{soluzioni-dark.pdf}),
allo stesso numero di esercizio. Il codice è anche in \texttt{codice/m11/}.
